\section{Method}

\subsection{Framework objective}

The method combines a forecasting benchmark, a verifier-gated candidate
proposal protocol, deterministic policy selection, and claim-boundary auditing.
The benchmark supplies evidence for weekly hospitalization-rate forecasting
under a common data-processing, fitting, and evaluation protocol. The
selection framework then turns candidates and compact evidence summaries into
auditable recommendations. The objective is threefold:
\begin{enumerate}[leftmargin=1.5em]
\item compare forecasting baselines, hand-specified mechanistic models,
structured-search candidates, and ablations under a shared evidence protocol;
\item decide which recommendations are credible when model families, proposer
orders, and evidence modes agree or disagree; and
\item prevent leakage and overclaiming by separating selection evidence,
post-hoc test evidence, verifier checks, and claim-boundary labels.
\end{enumerate}

All model selection and hyperparameter selection use only training,
validation, or rolling-validation evidence. Held-out test metrics are computed
only after the model family, structure, or baseline configuration has been
selected.

\subsection{Formal framework}

Let $y_{1:T}$ denote a univariate forecasting series. We use a chronological
split into training, validation, and held-out test windows,
$\mathcal{I}_{\mathrm{tr}}$, $\mathcal{I}_{\mathrm{val}}$, and
$\mathcal{I}_{\mathrm{test}}$. Candidate selection may use only evidence from
$\mathcal{I}_{\mathrm{tr}}$ and $\mathcal{I}_{\mathrm{val}}$, including
rolling-validation summaries when available. Held-out test evidence is
computed only after selection.

The candidate space is cross-family:
\[
\mathcal{C} =
\mathcal{C}_{\mathrm{forecast}} \cup
\mathcal{C}_{\mathrm{mech}} \cup
\mathcal{C}_{\mathrm{search}} \cup
\mathcal{C}_{\mathrm{ablation}} .
\]
A candidate record is represented as
\[
c = (\mathrm{family}, \mathrm{model\_name}, \mathrm{observation\_label},
\mathrm{delay\_label}, \mathrm{rationale}),
\]
with metadata such as proposer name, round index, expected failure mode, and
paired ablation target. At proposal round $r$, a proposer $\pi_\theta$
receives a no-leakage context $q_r$ and a remaining budget $B_r$, and emits at
most $b$ candidates:
\[
\pi_\theta(q_r, B_r) \rightarrow
\{c_{r,1}, \ldots, c_{r,b}\}.
\]
The proposer may be deterministic, random, failure-guided, mock API-assisted,
or a real API provider. In all cases, the output is advisory until it passes
the verifier.

The verifier is a conjunction of hard checks:
\[
V(c, q_r) =
V_{\mathrm{schema}} \wedge
V_{\mathrm{allowlist}} \wedge
V_{\mathrm{obs}} \wedge
V_{\mathrm{budget}} \wedge
V_{\mathrm{noLeak}} \wedge
V_{\mathrm{claim}} .
\]
These checks reject malformed records, duplicate identifiers, out-of-allowlist
families or models, incompatible observation or delay labels, budget
violations, absolute artifact paths, held-out test metrics in selection
evidence, and numerically flagged rows used to support positive claims.

For a verified candidate $c$, the selection evidence is
\[
E_{\mathrm{sel}}(c) =
(\mathrm{val\_mae}, \mathrm{rolling\_mae},
\mathrm{complexity}, \mathrm{numerical\_risk}, \mathrm{family}) .
\]
When validation MAE is unavailable in a compact replay table, the policy uses
rolling evidence and records that limitation. The held-out test evidence
$E_{\mathrm{test}}(c^\star)$ is attached only after a candidate $c^\star$ has
been selected and is labeled post-hoc descriptive.

The main policy, \texttt{pareto\_epsilon}, filters unsafe candidates and forms
an epsilon Pareto frontier over the selection objectives. Candidate $a$
epsilon-dominates candidate $b$ when it is no worse by more than $\epsilon$ on
each minimized objective and strictly better on at least one objective. Ties
are broken deterministically by numerical risk, rolling error, complexity, and
candidate identifier. For a budgeted ordering $c_1,\ldots,c_k$, define
\[
\operatorname{hit}_{\epsilon}(k) =
\begin{cases}
1, & \text{if any candidate among } c_1,\ldots,c_k
\text{ is within the top-epsilon reference set},\\
0, & \text{otherwise},
\end{cases}
\]
and let
\[
B_{\epsilon} = \min\{k : \operatorname{hit}_{\epsilon}(k)=1\}
\]
when such a $k$ exists. We also report deterministic weighted scoring and a
hard-veto decision tree as ablation policies.

Finally, the claim auditor maps evidence into allowed and rejected claims:
\[
A(\mathrm{evidence}) \rightarrow
(\mathrm{allowed\_claims}, \mathrm{rejected\_claims}, \mathrm{caveats}).
\]
This step records whether evidence supports group-specific signals, no single
global winner, or proposal-quality claims, and rejects claims such as
state-of-the-art forecasting, autonomous science, real-world mechanism
recovery, or using flagged rows as positive evidence.

\subsection{Verifier-gated selection layer}

The selection layer separates proposal, verification, evidence attachment,
policy selection, and claim-boundary auditing. Candidate records use a fixed
schema with a model family, model name, observation label, delay label,
proposal rationale, and metadata. The model-family space is cross-family:
forecasting baselines, hand-specified mechanistic baselines, structured search
models, ablation models, and ensembles.

The hard verifier rejects malformed records, duplicate identifiers, absolute
artifact paths, invalid family/model combinations, invalid observation or
delay labels, test metrics used as selection evidence, and numerically flagged
rows used to support positive claims. API-assisted proposal, when enabled, is
restricted to JSON output from an explicit allowlist and cannot generate model
code or execute candidate definitions. The default selection evaluations are
offline and deterministic.

The main policy is \texttt{pareto\_epsilon}, which filters unsafe candidates
and selects from an epsilon Pareto frontier over available selection evidence,
rolling-origin error, complexity, and numerical risk. We also report a
deterministic weighted-score policy as an ablation and a hard-veto decision
tree that prioritizes safety, simple baseline sufficiency, and
observation-label evidence. The claim-boundary auditor records whether the
evidence supports cautious statements such as group-specific signal and no
single global winner, and rejects statements such as global structured-search
superiority, state-of-the-art forecasting, or using flagged rows as positive
evidence.

\subsection{Constrained agentic proposer protocol}

The language-model component, when enabled, is a constrained agentic proposer
rather than a free-form scientist. The protocol is defined by four prompt-task
contracts, documented in \texttt{reports/agent\_prompt\_protocol.md}.
\begin{description}[leftmargin=1.5em]
\item[Task A: Initial candidate proposal] asks for a small, diverse,
allowlisted set of candidates for the next evaluation round.
\item[Task B: Evidence-aware refinement] provides verifier feedback,
accepted candidates, duplicate candidates, non-test rolling or validation
evidence, and remaining budget, then asks for refined candidates.
\item[Task C: Failure diagnosis or ablation request] converts non-test
failure evidence into an allowlisted ablation request.
\item[Task D: Claim-boundary summary] summarizes deterministic audit labels
without expanding the allowed claims.
\end{description}

Allowed prompt context includes the series name, candidate allowlist, budget,
round index, previous verifier rejection reasons, accepted candidates,
duplicate candidates, family diversity so far, observation labels already
tried, non-test rolling or validation evidence, numerical-risk warnings, and
remaining budget. Forbidden fields include held-out test MAE or RMSE, test
winners, test ranks, post-selection test comparisons, and paper
recommendations derived from test evidence. Provider outputs must be JSON or
provider-native structured outputs that match the shared schema; unknown model
names, executable-code fields, unsupported families, and claim-expanding
summaries are rejected before replay or execution.

\subsection{Data processing and series construction}

The raw FluSurv-NET CSV is read with \texttt{skiprows=2} to discard metadata
lines above the table. Rows with missing \texttt{WEEK} or \texttt{YEAR.1} are
removed, the remaining rows are sorted chronologically by year and week, and a
continuous integer time index is created. The primary series is the
\emph{Entire Network} / \emph{Overall} weekly hospitalization rate. The frozen
run also evaluates five age-stratified series:
\texttt{0--4 yr}, \texttt{5--17 yr}, \texttt{18--49 yr},
\texttt{50--64 yr}, and \texttt{$\geq$65 yr}.

\subsection{Forecasting baselines}

The discovery-ablation configuration adds opt-in non-epidemic forecasting
baselines. Persistence and rolling-mean baselines test whether recent observed
rates are already sufficient. For horizon $h$, the persistence baseline is
\[
\hat{y}_{T+h} = y_T,
\]
and a rolling mean with window $w$ is
\[
\hat{y}_{T+h} =
\frac{1}{w}\sum_{j=0}^{w-1} y_{T-j}.
\]
\texttt{arima\_auto\_small} chooses among a small fixed set of ARIMA orders by
validation MAE and then refits on train-plus-validation for test forecasting.
\texttt{lagged\_ridge} and \texttt{lagged\_gradient\_boosting} use lagged
observations, linear time, and seasonal sine/cosine features with recursive
multi-step forecasts. These baselines are opt-in; the historical default
benchmark remains the original six epidemic model families.

\subsection{Manual epidemic baselines}

The hand-specified epidemic baselines include deterministic SEIR,
probabilistic SEIR, hospitalized SEIHR, delayed-observation SEIR, and
fractional SEIR variants. All epidemic states are normalized, and fitting uses
explicit penalties for negative states and approximate mass-conservation
violations. The delayed-observation SEIR baseline selects its delay using
validation evidence only.

\subsection{Constrained structure discovery}

The discovery module searches over a bounded grammar rather than arbitrary
symbolic dynamics. A structure specification contains:
\begin{itemize}[leftmargin=1.5em]
\item a compartment template from
\{\texttt{SIR}, \texttt{SEIR}, \texttt{SEIRS}, \texttt{SEIHR}, \texttt{SEIAR}\},
\item a fractional on/off flag,
\item an observation map from
\{\texttt{I}, \texttt{H}, \texttt{I+H}, \texttt{delayed\_I}\}, and
\item a delay parameter in $\{0,1,2,3\}$ when the observation map is
\texttt{delayed\_I}.
\end{itemize}

The grammar templates are bounded SIR, SEIR, SEIRS, SEIHR, and SEIAR
state-space families. Let $x_t$ denote the latent state vector. Observations
are modeled as a scaled and noisy view of one allowed latent component or
component combination:
\[
y_t = \rho\, g_o(x_{t-d}) + \epsilon_t,
\]
where $g_o \in \{I_t, H_t, I_t + H_t, I_{t-d}\}$ when the required
compartments exist, $d$ is the allowed delay, and $\rho$ is a scale parameter.
The notation is intentionally schematic: the repository implements bounded
templates and validity checks rather than unrestricted equation generation.

Validity rules enforce that infection originates from $S$, infectious-like
states eventually reach recovery, no isolated compartments exist, and
observation semantics are compatible with the latent template. For example,
\texttt{obs=H} and \texttt{obs=I+H} require an $H$ compartment, while
\texttt{obs=delayed\_I} requires an $I$ compartment and a valid delay.

\subsection{Discovery score policies and ablations}

The constrained discovery model uses a stability-aware score that combines
multi-split validation error, rolling-origin stability, complexity penalties,
and the configured age prior. The frozen ablation package compares this search
against five same-grammar alternatives:
\[
\begin{aligned}
S_{\mathrm{stable}} &=
\overline{\mathrm{MAE}}_{\mathrm{multi\ split}} +
P_{\mathrm{multi\ split}} +
P_{\mathrm{stability}} +
P_{\mathrm{complexity}} +
P_{\mathrm{age}},\\
S_{\mathrm{validation}} &= \mathrm{MAE}_{\mathrm{val}},\\
S_{\mathrm{no\ stability}} &=
\overline{\mathrm{MAE}}_{\mathrm{multi\ split}} +
P_{\mathrm{multi\ split}} +
P_{\mathrm{complexity}} +
P_{\mathrm{age}}.
\end{aligned}
\]
\begin{itemize}[leftmargin=1.5em]
\item \texttt{random\_structure\_discovery}, which evaluates a seeded random
candidate budget;
\item \texttt{exhaustive\_structure\_discovery}, which evaluates the valid
grammar universe subject to an explicit candidate guard;
\item \texttt{validation\_only\_structure\_selection}, whose score is exactly
validation MAE;
\item \texttt{no\_observation\_search\_discovery}, which restricts the
candidate universe to \texttt{obs=I}; and
\item \texttt{no\_stability\_discovery}, which removes the rolling stability
penalty from the score while retaining other terms.
\end{itemize}

All discovery leaderboards include score-policy audit columns so that the
selected structure can be traced back to the evidence and penalties used.

\subsection{Synthetic structured recovery tasks}

The synthetic recovery tasks are generic structured time-series tests, not
domain simulations. They create a latent signal $z_t$, an auxiliary proxy
$u_t$, and observations with known labels:
\[
\begin{aligned}
\mathrm{direct}: \quad & y_t = z_t + \epsilon_t,\\
\mathrm{lagged}(d): \quad & y_t = z_{t-d} + \epsilon_t,\quad d\in\{1,2\},\\
\mathrm{mixture}: \quad & y_t = \alpha z_t + (1-\alpha)u_t + \epsilon_t,\\
\mathrm{proxy}: \quad & y_t = u_t + \epsilon_t.
\end{aligned}
\]
The task is to recover observation labels, delay labels, and useful candidate
families under fixed candidate budgets. These tasks validate software logic
for partial-observation candidate selection; they do not establish real-world
mechanism recovery.

\subsection{Post-hoc paired rolling comparison}

After model selection is fixed, the analysis aligns rolling-origin forecasts
by series, horizon, and target week. For each challenger and the constrained
discovery reference, it reports paired mean absolute-error differences,
bootstrap confidence intervals, and reference win rates. These paired
comparisons summarize post-hoc evidence and are not used for model selection.
